Los resultados experimentales se obtuvieron ejecutando ambos algoritmos sobre un conjunto fijo de 10 casos de prueba, previamente generados por un script de Python ubicado en \texttt{code/dynamic\_programming/scripts/input\_generator.py}. Cada caso contiene dos cadenas de texto $(s, t)$ con longitudes iguales, variando desde 5 hasta 160 caracteres, con diferencias estructuradas entre ellas.

Durante la ejecución, ambos programas escribieron las mediciones de tiempo (en milisegundos) y una estimación del uso de memoria (en KB) en archivos de texto ubicados en \texttt{data/measurements/a.txt}, los cuales pueden ser regenerados en cualquier entorno compatible simplemente ejecutando:

\begin{verbatim}
make              # desde code/brute_force o code/dynamic_programming
\end{verbatim}

Las gráficas fueron generadas a partir de estos archivos usando el script \texttt{plot\_generator.py}. El script lee las mediciones y construye dos gráficos: uno de tiempo de ejecución y otro de uso de memoria, guardándolos en \texttt{data/plots/}. Las figuras resultantes fueron exportadas como archivos PNG y agregadas al directorio \texttt{report/images/} para su inclusión en el informe.

\begin{figure}[H]
    \centering
    \begin{minipage}[t]{0.5\textwidth}
        \includegraphics[width=\textwidth]{images/tiempo_vs_tamanoDP.png}
        \caption*{(a) Tiempo vs Tamaño (DP)}
    \end{minipage}%
    \begin{minipage}[t]{0.5\textwidth}
        \includegraphics[width=\textwidth]{images/tiempo_vs_tamanoBF.png}
        \caption*{(b) Tiempo vs Tamaño (Fuerza Bruta)}
    \end{minipage}
    \caption{Tiempo de ejecución para ambos algoritmos}
    \label{fig:tiempos}
\end{figure}

\begin{figure}[H]
    \centering
    \begin{minipage}[t]{0.5\textwidth}
        \includegraphics[width=\textwidth]{images/memoria_vs_tamanoDP.png}
        \caption*{(c) Memoria vs Tamaño (DP)}
    \end{minipage}%
    \begin{minipage}[t]{0.5\textwidth}
        \includegraphics[width=\textwidth]{images/memoria_vs_tamanoBF.png}
        \caption*{(d) Memoria vs Tamaño (Fuerza Bruta)}
    \end{minipage}
    \caption{Uso de memoria estimado para ambos algoritmos}
    \label{fig:memoria}
\end{figure}

En la Figura \ref{fig:tiempos}, se aprecia que el algoritmo de programación dinámica escala de forma eficiente y estable, con tiempos de ejecución que se mantienen por debajo de 0.05 ms incluso para entradas grandes. En contraste, la fuerza bruta solo logra ejecutarse para los dos primeros casos, presentando un salto abrupto de tiempo entre entradas pequeñas, lo cual refleja su crecimiento exponencial.

La Figura \ref{fig:memoria} muestra que ambos algoritmos presentan un uso de memoria estimado constante, dado que se usó una estimación simplificada basada en estructuras de datos principales. Sin embargo, en la práctica, fuerza bruta requeriría mayor uso de pila por recursión profunda, lo cual podría generar errores de ejecución en casos grandes.

En conjunto, los resultados experimentales muestran que la programación dinámica no solo es teóricamente más eficiente, sino que también se comporta de forma significativamente mejor en la práctica, tanto en tiempo como en escalabilidad.